\documentclass[11pt]{article}
\usepackage[margin=0.75in]{geometry}
\usepackage{fontspec}
\usepackage{setspace}
\usepackage{hyperref}
\hypersetup{colorlinks=true, urlcolor=blue, linkcolor=black}
\pagestyle{empty}
\setlength{\parskip}{0.65em}
\setlength{\parindent}{0pt}
\setstretch{0.93}
\setmainfont{Times New Roman}

\begin{document}

\begin{flushleft}
Ke Cheng HSU (Clio) \\
Department of Economics, The Hong Kong University of Science and Technology \\
Clear Water Bay, Kowloon, Hong Kong $\cdot$ +852 6840 9753 \\
\href{mailto:khsu@connect.ust.hk}{khsu@connect.ust.hk} $\cdot$ \href{https://cliohsu1997.github.io/}{cliohsu1997.github.io}
\end{flushleft}

25 November 2026

\vspace{0.1em}
Faculty Search Committee \\
Department of Finance and Economics \\
Faculty of Business and Management \\
Beijing Normal--Hong Kong Baptist University (BNBU) \\
Zhuhai, China

\vspace{0.1em}
Dear Members of the Search Committee,

I am writing to apply for the Assistant Professor position in Applied Economics (Ref:\ FBM261003) in the Department of Finance and Economics. I learned of this opening from BNBU's Academic Positions page (\href{https://hro.bnbu.edu.cn/Job_Opportunities_/Job_Vacancy/Academic_Positions.htm}{hro.bnbu.edu.cn}). I am a Ph.D.\ candidate in Economics at the Hong Kong University of Science and Technology. I expect to graduate in June 2027 and can join for the February 2027 intake. My fields are Empirical Industrial Organization and Digital Economics. I have a B.Sc.\ in Mathematics and Economics (with a Minor in Actuarial Mathematics), and I am prepared to teach the AE Programme's mathematics-intensive courses---especially \textbf{Econometrics}, \textbf{Calculus}, and \textbf{Linear Algebra}---in English.

My job market paper, ``Designing Return Policy in E-Commerce: A Trade-off between Purchase and Return Efficiency,'' studies whether a return policy encourages efficient returns and purchases. A stricter policy curbs inefficient returns after delivery, but the lower return rate makes a higher product price more attractive for the seller and can exclude efficient purchases. Using transaction data from a Taobao apparel retailer, I estimate a structural model of purchase, return, and whether to buy return-shipping insurance under monopoly pricing across six categories. Automatic coverage of return shipping performs better when the margin, net scrap value, and uncertainty are high; optional coverage performs better when those attributes are low. A menu that pairs a discounted price with a medium restocking fee balances the two margins for most categories.

A second working paper, ``Product Design with Heterogeneous Search Costs'' (with Pak Hung Au), studies how firms choose product design when consumers differ in search costs. I also have two ongoing projects: with Kohei Kawaguchi and Kei Kawai, on resale price maintenance and unsold inventory waste using Japanese bookstore data; and with Ce Bian and Liang Zhao, on pass-through when products can be returned.

\textbf{Teaching is a core part of my fit for this post.} In 2026 I designed and lectured the \textbf{mathematics camp} for the MSc in Economics at HKUST (evaluation 4.53/5), covering the calculus and linear-algebra foundations that graduate students need for econometrics and theory. From 2020 to 2025 I was a teaching assistant for undergraduate macroeconomics (ECON2123), the economics of uncertainty and information (ECON4134 / ECON4999T), and Ph.D.\ microeconomic theory (ECON5210), with an average evaluation of 4.72/5. I also served as Teaching Assistant Coordinator in the Business School (2025--2026). I am ready to teach Econometrics, Calculus, Linear Algebra, applied microeconometrics, industrial organization, and digital economics / e-commerce from the undergraduate to the postgraduate level, and to supervise quantitative research projects in English.

I submit my CV and cover letter for the portal (Ref:\ FBM261003). Crystal, the secretary of the Department of Economics at HKUST, can arrange reference letters from Pak Hung Au and Kei Kawai if required. You can reach me at \href{mailto:khsu@connect.ust.hk}{khsu@connect.ust.hk} or +852 6840 9753.

\vspace{0.2em}
Sincerely,

\vspace{0.2em}
Ke Cheng HSU \\
Ph.D.\ Candidate, Department of Economics, HKUST

\end{document}
